\scanpage{63}
\mspage{27}
\deleted{c) Si $(z_\alpha)_{\alpha\in I}$ est l'ens des éléments de
$\Sigma_{<x}=\uline{b}^{-1}(\{x\})$ (où $\uline{b}:\Sigma^{<}(1)\longrightarrow\Sigma$), les
\unsure{les}}

a) Pour tt $x\in\Sigma$, l'ens $\Sigma_{<x}$ \deleted{(fini) est}\inserted{est tot${}^t$}
\deleted{$\simeq\uline{b}^{-1}(\{x\})$} ordonné

b) Si $\Sigma^{<}(1)$ est le graphe de la relation
$x<y$, et $\uline{s},\uline{b}:\Sigma^{<}(1)\rightrightarrows\Sigma$ sont les applications
source et but, \deleted{\unsure{---}} l'application $\uline{b}$ est une immer-
sion \deleted{bordante} \uline{propre} (donc les $\Sigma_{<x}=\uline{b}^{-1}(\{x\})$
sont finis)

c) Si $(z_\alpha)_{\alpha\in I}$ sont les éléments de $\Sigma_{<x}=b^{-1}(\{x\})$,
alors \deleted{les}\inserted{la famille} des germes \inserted{en $x$} de sous-espaces \deleted{bordants}\inserted{de $\Sigma$}
images par $\uline{b}$ des \inserted{germes} $(\Sigma^{<}(1),z_\alpha)$, \deleted{sont}\inserted{une famille}
\deleted{transversaux}\inserted{transversale}
de sous-espaces \uline{bordants}.

d) $\uline{s}:\Sigma^{<}(1)\longrightarrow\Sigma^{\unsure{*}}$ est \uline{transversal} p.r. à l'immersion
bordante $\uline{b}:\Sigma^{<}(1)\longrightarrow\Sigma^{\unsure{*}}$
\deleted{\unsure{la diag}}
\margin{ou sous forme d'une \unsure{récollection|récollement}
pour les $z_\alpha\in\Sigma^{<}(1)$ au dessus d'un même
$x\in\Sigma^*$\ldots}

$1^{\circ})$ Cas équisingulier : $\uline{s}:\Sigma^{<}(1)\longrightarrow\Sigma$ fibrant \inserted{(propre)}

$2^{\circ})$ Cas non singulier : $\Sigma$ est une variété à bord,
dont le bord est $\uline{b}(\Sigma^{*}(1))$.

$3^{\circ})$ Cas équisingulier-non singulier : $\uline{s}:\Sigma^{<}(1)\longrightarrow\Sigma$ est
fibrant, \inserted{propre} à fibres des variétés à bord, et le
fibré des bords est l'image de $\Sigma^{*}(2)$ par
le morphisme $\Sigma^{*}(2)\xrightarrow{\uline{s'}}\Sigma^{*}(1)$ \deleted{\unsure{avec}}\inserted{donné}
par $(x\prec y\prec z)\longrightarrow(x\prec z)$.
\margin{c'est sans doute conséquence de la conjonction des deux conditions $1^{\circ}),2^{\circ})$}
\mspage{62}
